% !TEX root = ../main.tex
\corrigesChapitre{18}{Modèle du gaz parfait}
\corrigesSouples

\begin{blocLivre}{cor-18-1}
\begin{corrige}{1}{Interpréter à l'échelle microscopique}
\demarche{Masse volumique : nombre d'entités par unité de volume ; pression : chocs sur les parois ; température :
agitation.}
\begin{questions}
\item Le volume, la masse volumique (nombre et masse des molécules par unité de volume), la pression (\res{\text{chocs}} des
molécules sur les parois) et la température (\res{\text{agitation}} des molécules).
\item Le même nombre de molécules occupe un volume plus petit : la masse volumique \res{\text{augmente}}. Les molécules,
plus serrées, frappent plus souvent les parois : la pression \res{\text{augmente}}.
\item La masse volumique \res{\text{ne change pas}} (même nombre de molécules, même volume). Les molécules, plus agitées,
frappent les parois plus souvent et plus violemment : la pression \res{\text{augmente}}.
\item \qty{293,15}{K}, \qty{255,15}{K} et \qty{373,15}{K}. Une température de \qty{-300}{\degreeCelsius} correspondrait à
$T<0$ : elle serait inférieure au \res{\text{zéro absolu}}, ce qui est impossible.
\end{questions}
\end{corrige}
\end{blocLivre}
\begin{blocLivre}{cor-18-2}
\begin{corrige}{2}{L'air d'une chambre}
\demarche{$V_\mathrm{m}=R\,T/p$ et $n=p\,V/(R\,T)$, avec $V=4{,}0\times5{,}0\times2{,}5=\qty{50}{m^3}$.}
\begin{questions}
\item $V_\mathrm{m}=\dfrac{8{,}314\times273{,}15}{\num{1.013e5}}=\res{\qty{22,4}{L.mol^{-1}}}$ et, à \qty{293,15}{K},
\res{\qty{24,1}{L.mol^{-1}}}.
\item $n=\dfrac{\num{1.013e5}\times50}{8{,}314\times293{,}15}=\res{\qty{2,08e3}{mol}}$.
\item $m=n\,M=\res{\qty{60,3}{kg}}$ : l'air de la chambre pèse autant qu'un adulte.
\item À pression et volume constants, $n$ est inversement proportionnelle à $T$ : $n'=n\times\dfrac{293{,}15}{298{,}15}=\qty{2,04e3}{mol}$.
Il sort environ \res{\qty{35}{mol}} d'air, soit \res{\qty{1,0}{kg}}.
\end{questions}
\end{corrige}
\end{blocLivre}
\begin{blocLivre}{cor-18-3}
\begin{corrige}{3}{La seringue bouchée}
\demarche{Quantité d'air et température constantes : $p\,V$ est constant.}
\begin{questions}
\item La \res{\text{loi de Mariotte}}, $p\,V=\text{constante}$ : la seringue est bouchée ($n$ constant) et la température ne
change pas.
\item $p=\dfrac{1013\times50}{20}=\res{\qty{2,53e3}{hPa}}$.
\item $F=(p-p_\mathrm{atm})\,S=(\num{2.533e5}-\num{1.013e5})\times\num{2.5e-4}=\res{\qty{38}{N}}$.
\item $p=\dfrac{1013\times50}{80}=\res{\qty{633}{hPa}}$.
\end{questions}
\end{corrige}
\end{blocLivre}
\begin{blocLivre}{cor-18-4}
\begin{corrige}{4}{Le pneu qui chauffe}
\demarche{Volume et quantité d'air constants : $p/T$ est constant, avec la pression absolue et $T$ en kelvins.}
\begin{questions}
\item $p_1=2{,}30+1{,}00=\res{\qty{3,30}{bar}}$.
\item $p_2=p_1\times\dfrac{323{,}15}{283{,}15}=\res{\qty{3,77}{bar}}$ : le manomètre indique \res{\qty{2,77}{bar}}.
\item $\dfrac{323{,}15}{283{,}15}=1{,}14$ : une hausse de \res{\qty{14}{\percent}}.
\item À chaud, la pression lue est \res{\text{plus élevée}} : régler les pneus à \qty{2,30}{bar} à chaud les laisserait
sous-gonflés une fois refroidis.
\end{questions}
\end{corrige}
\end{blocLivre}
\begin{blocLivre}{cor-18-5}
\begin{corrige}{5}{La masse volumique d'un gaz}
\demarche{$\rho=m/V$ avec $m=n\,M$ et $n=p\,V/(R\,T)$.}
\begin{questions}
\item $\rho=\dfrac{n\,M}{V}=\dfrac{p\,V}{R\,T}\times\dfrac{M}{V}=\res{\dfrac{p\,M}{R\,T}}$.
\item Air : \res{\qty{1,21}{kg.m^{-3}}} ; hélium : \res{\qty{0,166}{kg.m^{-3}}} ; dioxyde de carbone : \res{\qty{1,83}{kg.m^{-3}}}.
\item L'hélium est sept fois moins dense que l'air : la poussée d'Archimède exercée par l'air sur le ballon
\res{\text{dépasse son poids}}.
\item Le dioxyde de carbone est \res{\text{1,52 fois plus dense}} que l'air : il s'accumule au fond de la cuve et chasse
l'air respirable.
\end{questions}
\end{corrige}
\end{blocLivre}
\begin{blocLivre}{cor-18-6}
\begin{corrige}{6}{La bouteille de plongée}
\demarche{Équation d'état pour la bouteille ; à température donnée, $p\,V$ se conserve quand l'air se détend.}
\begin{questions}
\item $n=\dfrac{200\times10^5\times\num{12e-3}}{8{,}314\times293{,}15}=\res{\qty{98,5}{mol}}$.
\item $V=12\times\dfrac{200}{1}=\res{\qty{2400}{L}}$.
\item $V$ et $n$ constants : $p=200\times\dfrac{285{,}15}{293{,}15}=\res{\qty{194,5}{bar}}$.
\item L'air utilisable correspond à $12\times(194{,}5-50)=\qty{1734}{L}$ sous \qty{1}{bar}. À \qty{3,0}{bar}, chaque minute
consomme l'équivalent de $20\times3{,}0=\qty{60}{L}$ sous \qty{1}{bar} : autonomie $\dfrac{1734}{60}\approx\res{\qty{29}{min}}$.
\end{questions}
\end{corrige}
\end{blocLivre}
\begin{blocLivre}{cor-18-7}
\begin{corrige}{7}{La bulle du plongeur}
\demarche{Quantité d'air constante : $\dfrac{p_1V_1}{T_1}=\dfrac{p_2V_2}{T_2}$.}
\begin{questions}
\item $p_1=\num{1.013e5}+1025\times9{,}81\times20=\res{\qty{3,02e5}{Pa}}$.
\item $V_2=V_1\times\dfrac{p_1}{p_2}\times\dfrac{T_2}{T_1}=1{,}0\times\dfrac{3{,}024}{1{,}013}\times\dfrac{291{,}15}{285{,}15}=\res{\qty{3,0}{cm^3}}$.
\item À température constante, on obtiendrait \qty{2,99}{cm^3} : le réchauffement n'augmente le volume que de
\res{\qty{2}{\percent}} environ ; l'essentiel vient de la baisse de pression.
\item Le volume serait multiplié par 3, soit \res{\qty{18}{L}} : les poumons seraient gravement lésés (surpression
pulmonaire). Règle : \res{\text{toujours expirer}} pendant la remontée.
\end{questions}
\end{corrige}
\end{blocLivre}
\begin{blocLivre}{cor-18-8}
\begin{corrige}{8}{Tester l'équation d'état en Python}
\demarche{Conversions : $\qty{1}{hPa}=\qty{e2}{Pa}$ et $\qty{1}{mL}=\qty{e-6}{m^3}$ ; $n=p\,V/(R\,T)$.}
\begin{questions}
\item \code{1e2} convertit les hectopascals en pascals et \code{1e-6} les millilitres en mètres cubes. Le programme affiche
le produit \res{p\,V} (en \unit{hPa.mL}) et la \res{\text{quantité de matière}} $n$ (en millimoles).
\item $p\,V$ reste égal à \qty{60800}{hPa.mL} à 0,2~\% près : il est \res{\text{constant}} aux incertitudes de mesure près
(lecture du volume à \qty{0,5}{mL} près, soit 1 à 3~\%). La loi de Mariotte est vérifiée.
\item \code{inv = [1 / x for x in V]}, \code{plt.plot(inv, p, "o")}, puis \code{plt.show()} : des points alignés sur une
\res{\text{droite passant par l'origine}}, de coefficient directeur $p\,V=n\,R\,T\approx\qty{60800}{hPa.mL}=\res{\qty{6,08}{J}}$.
\item $V$ et $n$ constants : $p=3035\times\dfrac{308{,}15}{293{,}0}=\res{\qty{3,19e3}{hPa}}$.
\end{questions}
\schema[corrige]{rev-18-ex08-cor}
\legende{Le tracé de $p$ en fonction de $1/V$ permet de tester la proportionnalité.}
\end{corrige}
\end{blocLivre}
\begin{blocLivre}{cor-18-9}
\begin{corrige}{9}{À la recherche du zéro absolu}
\demarche{Droite $p=a\,\theta+b$ ; la pression s'annule pour $\theta=-b/a$.}
\begin{questions}
\item Les points sont \res{\text{alignés}} : $p$ est une fonction affine de $\theta$, croissante.
\item $a=\dfrac{1130-881}{90-10}=\res{\qty{3,11}{hPa.\degreeCelsius^{-1}}}$ et $b=881-3{,}11\times10=\res{\qty{850}{hPa}}$.
\item $p=0$ pour $\theta=-\dfrac{b}{a}=\res{\qty{-273}{\degreeCelsius}}$, en accord avec le zéro absolu, \qty{-273,15}{\degreeCelsius}.
\item $p=a\,(\theta+273)=a\,T$ : la pression est \res{\text{proportionnelle à }T}, comme l'indique $p=\dfrac{n\,R}{V}\,T$ à $n$ et
$V$ constants.
\end{questions}
\schema[corrige]{rev-18-ex09-cor}
\legende{La droite est prolongée mathématiquement ; l’air réel ne demeure pas gazeux sur tout ce domaine.}
\end{corrige}
\end{blocLivre}
\begin{blocLivre}{cor-18-10}
\begin{corrige}{10}{La bouteille écrasée}
\demarche{À $n$ et $V$ constants, $p/T$ est constant.}
\begin{questions}
\item $p=1013\times\dfrac{255{,}15}{298{,}15}=\res{\qty{867}{hPa}}$.
\item $\Delta p=1013-867=\res{\qty{146}{hPa}}$.
\item $F=\Delta p\times S=\num{14600}\times\num{1.0e-2}=\res{\qty{146}{N}}$, vers l'intérieur : c'est le poids d'une masse
d'environ \res{\qty{15}{kg}}.
\item À $n$ et $p$ constants, $V$ est proportionnel à $T$ : $V=1{,}5\times\dfrac{255{,}15}{298{,}15}=\res{\qty{1,28}{L}}$.
\end{questions}
\end{corrige}
\end{blocLivre}
\begin{blocLivre}{cor-18-11}
\begin{corrige}{11}{Les limites du modèle}
\demarche{Écart relatif $\dfrac{V_\mathrm{m}-V_\mathrm{m,parfait}}{V_\mathrm{m,parfait}}$ ; un gaz dense ou liquéfié échappe au modèle.}
\begin{questions}
\item $V_\mathrm{m}=\dfrac{8{,}314\times273{,}15}{\num{1.01325e5}}=\res{\qty{22,41}{L.mol^{-1}}}$.
\item \ce{N2} : \qty{-0,04}{\percent} ; \ce{O2} : \qty{-0,09}{\percent} ; \ce{CO2} : \qty{-0,7}{\percent} ; \ce{NH3} : \qty{-1,4}{\percent}.
\ce{CO2} et \ce{NH3} s'écartent le plus : leurs molécules, plus grosses ou polaires, exercent entre elles des
\res{\text{interactions attractives}} plus fortes, qui rapprochent les molécules.
\item $n=\dfrac{2000}{44{,}0}=\qty{45,5}{mol}$ et $p=\dfrac{n\,R\,T}{V}=\res{\qty{369}{bar}}$.
\item Au-delà de \qty{57}{bar}, le dioxyde de carbone \res{\text{se liquéfie}} : l'extincteur contient surtout du liquide,
surmonté de gaz à \qty{57}{bar}. Le modèle du gaz parfait ne s'applique plus.
\end{questions}
\end{corrige}
\end{blocLivre}
\begin{blocLivre}{cor-18-12}
\begin{corrige}{12}{Le coussin gonflable de sécurité}
\demarche{Équation d'état pour le diazote, puis proportions de l'équation : $n(\ce{NaN3})=\frac23\,n(\ce{N2})$.}
\begin{questions}
\item $n(\ce{N2})=\dfrac{\num{1.30e5}\times\num{0.060}}{8{,}314\times323{,}15}=\res{\qty{2,90}{mol}}$.
\item $n(\ce{NaN3})=\dfrac23\times2{,}90=\res{\qty{1,94}{mol}}$.
\item $m=1{,}94\times65{,}0=\res{\qty{126}{g}}$.
\item $p=1{,}30\times\dfrac{293{,}15}{323{,}15}=\res{\qty{1,18}{bar}}$. Un coussin resté tendu ferait rebondir l'occupant : les
évents le \res{\text{dégonflent}} pour amortir le choc.
\end{questions}
\end{corrige}
\end{blocLivre}
\begin{blocLivre}{cor-18-13}
\begin{corrige}{13}{Le ballon-sonde éclate}
\demarche{Quantité d'hélium constante : $V=\dfrac{n\,R\,T}{p}$ à chaque altitude ; éclatement pour $V=\dfrac{\pi\times8{,}0^3}{6}$.}
\begin{questions}
\item $n=\dfrac{\num{1.013e5}\times3{,}0}{8{,}314\times288{,}15}=\res{\qty{127}{mol}}$.
\item $V=\dfrac{127\times8{,}314\times226{,}5}{1197}=\res{\qty{200}{m^3}}$, soit $d=\left(\dfrac{6V}{\pi}\right)^{1/3}=\res{\qty{7,3}{m}}$,
contre \qty{1,79}{m} au sol.
\item L'enveloppe éclate pour $V=\qty{268}{m^3}$. À \qty{30}{km}, $V=\qty{200}{m^3}$ ; à \qty{35}{km}, $V=\qty{434}{m^3}$ : le ballon
éclate \res{\text{entre 30 et 35~km}} d'altitude.
\item L'hélium est sous très \res{\text{faible pression}} et très loin de sa température de liquéfaction (\qty{4}{K}) : il
est très peu dense, ses atomes n'interagissent presque pas.
\end{questions}
\end{corrige}
\end{blocLivre}
\begin{blocLivre}{cor-18-14}
\begin{corrige}{14}{Respirer au sommet de l'Everest}
\demarche{$\rho=\dfrac{p\,M}{R\,T}$ ; dans un volume $V$, $n(\ce{O2})=0{,}21\,\dfrac{p\,V}{R\,T}$.}
\begin{questions}
\item Niveau de la mer : $\rho=\dfrac{\num{1.013e5}\times\num{0.0290}}{8{,}314\times288{,}15}=\res{\qty{1,23}{kg.m^{-3}}}$ ; sommet :
\res{\qty{0,474}{kg.m^{-3}}}.
\item Niveau de la mer : \res{\qty{8,9}{mmol}} de dioxygène par litre ; sommet : \res{\qty{3,4}{mmol}}, soit 39~\% seulement.
\item Dans les poumons, l'air est à \qty{37}{\degreeCelsius} dans les deux cas : à volume égal, $n$ est proportionnelle à $p$,
d'où un rapport $\dfrac{337}{1013}=\res{0{,}33}$, le tiers environ.
\item Le sang reçoit trois fois moins de dioxygène à chaque inspiration : malgré une respiration accélérée, l'organisme
\res{\text{manque de dioxygène}} (hypoxie), d'où l'usage de bouteilles.
\end{questions}
\end{corrige}
\end{blocLivre}
\begin{blocLivre}{cor-18-15}
\begin{corrige}{15}{Pourquoi un moteur Diesel n'a pas de bougie}
\demarche{Quantité d'air constante : $\dfrac{p_1V_1}{T_1}=\dfrac{p_2V_2}{T_2}$ avec $V_2=V_1/16$.}
\begin{questions}
\item $n=\dfrac{\num{1.0e5}\times\num{500e-6}}{8{,}314\times303{,}15}=\res{\qty{0,0198}{mol}}$, soit $m=\res{\qty{0,575}{g}}$.
\item À température constante, $p_2=16\,p_1=\res{\qty{16}{bar}}$.
\item $T_2=T_1\times\dfrac{p_2}{p_1}\times\dfrac{V_2}{V_1}=303{,}15\times\dfrac{45}{16}=\res{\qty{853}{K}}$, soit \res{\qty{579}{\degreeCelsius}},
bien au-dessus de la température d'auto-inflammation du gazole : il s'enflamme dès son injection.
\item Un mélange d'air et d'essence trop comprimé \res{\text{s'enflammerait spontanément}} avant le moment voulu (cliquetis),
ce qui endommage le moteur : la combustion doit être déclenchée par la bougie.
\end{questions}
\end{corrige}
\end{blocLivre}
